\documentclass{article}


\usepackage[utf8]{inputenc}   % MUST come before biblatex
\usepackage[T1]{fontenc}
\usepackage{csquotes}         % recommended for biblatex

\usepackage[
    backend=biber,
    style=alphabetic, % or whatever you want
]{biblatex}

% \usepackage[round]{natbib}
% \renewcommand{\bibsection}{\subsubsection*{\bibname}}

\usepackage[utf8]{inputenc} % allow utf-8 input
\usepackage[T1]{fontenc}    % use 8-bit T1 fonts
\usepackage{hyperref}       % hyperlinks
\usepackage{url}            % simple URL typesetting
\usepackage{booktabs}       % professional-quality tables
\usepackage{amsfonts}       % blackboard math symbols
\usepackage{nicefrac}       % compact symbols for 1/2, etc.
% \usepackage{microtype}      % microtypography
\usepackage{xcolor}         % colors
\usepackage[normalem]{ulem} % for \sout (strikethrough)

% \usepackage[bookmarks,colorlinks,breaklinks]{hyperref}
\hypersetup{urlcolor=blue, colorlinks=true, citecolor=green!50!black, linkcolor=blue}
\usepackage{multirow}
\usepackage[T1]{fontenc} % For polish hook https://tex.stackexchange.com/questions/75396/how-to-use-polishhook-symbol
\usepackage{amsmath}
\usepackage{amsfonts}
\usepackage{amssymb}
\usepackage{amsthm}
\usepackage{thmtools}
\usepackage{pgffor}
\usepackage{xcolor}
% \usepackage[lined,boxed,ruled,norelsize,algo2e,linesnumbered,noend]{algorithm2e}
\usepackage{cleveref}
\usepackage{graphicx}
\usepackage{geometry}
% \geometry{verbose,tmargin=1in,bmargin=1in,lmargin=1in,rmargin=1in}
\usepackage{enumitem}


% \newtheoremstyle{mystyle}
%   {12pt} % Space above
%   {24pt} % Space below
%   {} % Body font
%   {} % Indent amount
%   {\bfseries} % Theorem head font
%   {.} % Punctuation after theorem head
%   {.5em} % Space after theorem head
%   {} % Theorem head spec (can be left empty, meaning `normal')

% \theoremstyle{mystyle}


\newtheorem{assumption}{Assumption}
\newtheorem{definition}{Definition}
\newtheorem{theorem}{Theorem}
\newtheorem{lemma}{Lemma}
\newtheorem{proposition}{Proposition}
\newtheorem{corollary}{Corollary}

\usepackage{thmtools}

% \declaretheorem[refname={Theorem,Theorems},style=mystyle, Refname={Theorem,Theorems}]{theorem}
% \declaretheorem[numberlike=theorem, style=mystyle]{lemma}
% \declaretheorem[numberlike=theorem, style=mystyle]{invariant}
% \declaretheorem[numberlike=theorem, style=mystyle]{corollary}
% \declaretheorem[numberlike=theorem, style=mystyle]{definition}
% \declaretheorem[numberlike=theorem, style=mystyle]{claim}
% \declaretheorem[numberlike=theorem, style=mystyle]{fact}
% \declaretheorem[numberlike=theorem, style=mystyle]{assumption}
% \declaretheorem[numbered=no]{remark}

% number sets
\newcommand{\Real}{\mathbb{R}}
\newcommand{\Nat}{\mathbb{N}}
\newcommand{\F}{\mathbb{F}}
\newcommand{\Z}{\mathbb{Z}}
\renewcommand{\S}{\mathbb{S}}%

% probability
\renewcommand{\P}{\mathbb{P}}
\newcommand{\E}{\mathbb{E}}

% redfine some classical notation
\renewcommand{\tilde}{\widetilde}
\renewcommand{\hat}{\widehat}
\renewcommand{\bar}{\overline}

% operators
\newcommand{\median}{\operatorname{median}}

\DeclareMathOperator*{\argmin}{argmin}
\DeclareMathOperator*{\argmax}{argmax}
\newcommand{\poly}{\operatorname{poly}}
\newcommand{\polylog}{\operatorname{polylog}}
\newcommand{\mdiag}{\mathbf{Diag}} % diagonal matrix with vector on diag
\newcommand{\sign}{\operatorname{sign}}%
\newcommand{\nnz}{\operatorname{nnz}}%
\newcommand{\tail}{\operatorname{tail}}%
\newcommand{\dist}{\operatorname{dist}}%

% matrices
% \mX is \mathbf{X}
\foreach \x in {A,...,Z}{%
	\expandafter\xdef\csname m\x\endcsname{\noexpand\mathbf{\x}}
}
% \omX is \overline{\mathbf{X}}
\foreach \x in {A,...,Z}{%
	\expandafter\xdef\csname om\x\endcsname{\noexpand\overline{\noexpand\mathbf{\x}}}
}

\foreach \x in {A,...,Z}{%
	\expandafter\xdef\csname c\x\endcsname{\noexpand\mathcal{\x}}
}




\newcommand{\ob}{\overline{b}}
\newcommand{\ox}{\overline{x}}
\newcommand{\os}{\overline{s}}


\usepackage{mathtools}


% Reduce paragraph spacing https://tex.stackexchange.com/questions/4891/how-do-i-control-the-spacing-above-a-new-paragraph
\makeatletter
\renewcommand{\paragraph}{%
	\@startsection{paragraph}{4}%
	{\z@}{1.25ex \@plus 1ex \@minus .2ex}{-1em}%
	{\normalfont\normalsize\bfseries}%
}
\makeatother

\def\norm#1{\left\| #1 \right\|}



% \usepackage[
%     sorting=nyt,
% 	style=alphabetic,
%     maxalphanames=3,
%     minalphanames=3,
%     maxcitenames=3,
%     minnames=1,
%     maxbibnames=20,
% 	backref=true,
% 	doi=true,
% 	url=true,
% 	backend=biber
% ]{biblatex}
\usepackage{tcolorbox}


\newcommand{\thmbox}[1]{%
  \begin{tcolorbox}[
      colback = blue!5!white,
      colframe = blue!75!black,
      width          = \dimexpr\textwidth + 1cm\relax,
      % width    = \textwidth,   % full text width
      % left     = 5pt,          % no extra padding on the left
      % right    = 5pt,           % no extra padding on the right
      enlarge left by  = -0.5cm,
      enlarge right by = -0.5cm
  ]
    #1
  \end{tcolorbox}%
}

\newcommand{\mainthmbox}[1]{%
  \begin{tcolorbox}[colback=blue!5!white,colframe=blue!75!black,boxrule=3pt,
    ]
    #1
  \end{tcolorbox}%
}


\newcommand{\defbox}[1]{%
  \begin{tcolorbox}[colback=orange!5!white,colframe=orange!75!black, 
    width          = \dimexpr\textwidth + 1cm\relax,
    % width    = \textwidth,   % full text width
    % left     = 5pt,          % no extra padding on the left
    % right    = 5pt,           % no extra padding on the right
    enlarge left by  = -0.5cm,
    enlarge right by = -0.5cm]
    #1
  \end{tcolorbox}%
}

\newcommand{\blockbox}[1]{%
  \begin{tcolorbox}[colback=yellow!5!white,colframe=yellow!75!black]
    #1
  \end{tcolorbox}%
}

\newcommand{\shortminus}{\scalebox{0.5}[1.0]{\( - \)}}

% \usepackage{mathabx}  

\usepackage{cancel}
% \usepackage{arydshln}

\usepackage{geometry}
 \geometry{
        a4paper,
        total={170mm,257mm},
        left=20mm,
        right=10mm,
        top=20mm,
        }
        
% \usepackage{graphicx} % Required for inserting images

\usepackage{tikz}

\usetikzlibrary{arrows.meta,
                positioning,
                shadows}
                
\usetikzlibrary{shapes,snakes}
\usepackage{bm}
\usepackage{xspace}

\usepackage{cancel}


\setlist[itemize]{topsep=1pt, itemsep=1pt, parsep=1pt, leftmargin=*}

\usepackage{subcaption}
% \usepackage{graphicx}


\newcommand{\M}{\mathcal{M}}
\newcommand{\measure}{\mu}


% \newcommand{\K}{\mathbf{T}}
% \newcommand{\TW}{\mathbf{T}_W}
% \newcommand{\Ktilde}{\Tilde{\K}}
% \newcommand{\Khat}{\hat{\K}}

\newcommand{\K}{\bm{\mL}}
\newcommand{\TW}{\Tilde{\bm{\mL}}}
\newcommand{\Ktilde}{\Tilde{\K}}
\newcommand{\Khat}{\hat{\K}}

\newcommand{\RW}[1]{\mathbf{R}_W^{(#1)}}
\newcommand{\R}[1]{\mathbf{R}^{(#1)}}
\newcommand{\Rtilde}[1]{\Tilde{\mathbf R}^{(#1)}}
\newcommand{\Rhat}[1]{\hat{\mathbf R}^{(#1)}}

% \newcommand{}{}
\newcommand{\rc}[1]{\tilde r^{(#1)}}
\newcommand{\rt}[1]{r^{(#1)}}
\newcommand{\rtilde}[1]{\Tilde{r}^{(#1)}}
\newcommand{\rhat}[1]{\hat{r}^{(#1)}}

\newcommand{\rrho}{\mathbf{\tilde r}}
\newcommand{\rr}{\mathbf{r}}
\newcommand{\rrtilde}{\Tilde{\rr}}
\newcommand{\rrhat}{\hat{\rr}}


\newcommand{\PE}{\mathsf{p}}
\newcommand{\PEd}{\mathsf{p}}
\newcommand{\PEtilde}{\Tilde{\PE}}
\newcommand{\PEhat}{\hat{\PE}}

% creates \va to \vz vectors
\foreach \x in {a,...,z}{%
	\expandafter\xdef\csname v\x\endcsname{{\noexpand\boldsymbol{\noexpand\mathrm{\x}}}}
}


\newcommand{\signal}{\psi}
\newcommand{\opnorm}{op}
% \newcommand{\vv}{\mathrm{v}}
\newcommand{\xx}{x}
\newcommand{\yy}{\vy}
\newcommand{\zz}{\vz}
\newcommand{\uu}{\vu}


\newcommand{\NN}{N}

\newcommand{\Lip}[1]{\left\|#1\right\|_{\mathrm{Lip}}}
\newcommand{\LocLip}[2]{\left\|#1\right\|_{\mathrm{Lip}\left(#2\right)}}
\newcommand{\LLip}[1]{L_{#1}}
\newcommand{\Lipf}[1]{l(#1)}



\newcommand{\msg}{\mathsf{v}}
\newcommand{\upd}{\mathsf{W}}
\newcommand{\aggout}{\alpha_1}
\newcommand{\aggin}{\alpha_2}

\newcommand{\vpi}{\mathsf{\PE}}
\newcommand{\CM}{\mathsf{M}}
\renewcommand{\next}{\bm{\mathsf{next}}}
\newcommand{\att}{\mathsf{logit}}
\newcommand{\Att}{\mathsf{Att}}
\newcommand{\logit}{\mathsf{Logit}}
\newcommand{\msgc}{\tilde M}
\newcommand{\Msg}{\mathsf{V}}
\newcommand{\msghat}{\hat V}
\newcommand{\outhat}{\hat \out}
\newcommand{\hst}{\mathsf{h}}
\newcommand{\pool}{\mathsf{pool}}
\newcommand{\nn}{\mathsf{NN}}

\newcommand{\event}{\mathcal{E}}

\newcommand{\sparsity}{\gamma}

% eigenfunctions, eigenvalues and eigenvectors
\newcommand{\eigf}{v}
\newcommand{\eigv}{\vv}

\newcommand{\proj}{\Pi}

\newcommand{\find}{\operatorname{closest}}
\newcommand{\mLambda}{\bm{\Lambda}}
\newcommand{\gap}{\mathrm{gap}}


\newcommand{\downsampling}[1]{{D}_X #1 }
\newcommand{\interpolation}[1]{{U}_\chi #1	}

\newcommand{\Adj}{\operatorname{Adj}}
\newcommand{\Deg}{\operatorname{Deg}}
\newcommand{\Lap}{\operatorname{Lap}}

% \newcommand{\adj}{\mA}

\newcommand{\id}{\operatorname{id}}

\newcommand{\Q}{\mathrm{Q}}

\newcommand{\rankd}{\operatorname{rank}_d}
\newcommand{\absrank}{\operatorname{absrank}}
\newcommand{\rank}{\operatorname{rank}}

\newcommand{\hs}{{\scriptscriptstyle \mathrm{HS}}}
\newcommand{\op}{{\scriptscriptstyle \mathrm{op}}}

\newcommand{\pan}[1]{\textcolor{blue}{Pan:~#1}}
\newcommand{\todo}[1]{\textcolor{red}{TODO:~#1}}
\newcommand{\todoish}[1]{\textcolor{red}{#1}}
\newcommand{\nastia}[1]{\textcolor{orange}{Nastia:~#1}}
\newcommand{\mSigma}{\bm{\Sigma}}


\usepackage{amsmath}
\usepackage{xfrac}

% \foreach \x in {A,...,Z}{%
% 	\expandafter\xdef\csname op\x\endcsname{{\noexpand\boldsymbol{\noexpand{\mathrm{\x}}}}}
% }

\newcommand{\permLLip}[1]{\tilde L_{#1}}

\newcommand{\sampledM}{\mathrm{S}}

\newcommand{\rvx}{X}
\newcommand{\rvz}{Z}


\newcommand{\kernel}{\kappa}
% \newcommand{\mtrx}{\mK}
\newcommand{\ssX}{\bm{X}}
\newcommand{\sampledset}{\bm{X}}

\newcommand{\shortequals}{\mathrel{\mathop=}\limits^{\textstyle}}

\newcommand{\GNN}{\mathrm{GNN}}
\newcommand{\CGNN}{\mathrm{CGNN}}
\newcommand{\MPNN}{\mathrm{MPNN}}
% \newcommand{\GAT}{\mathrm{GAT}}
\newcommand{\GAT}{\Theta}
\newcommand{\CMPNN}{\mathrm{CMPNN}}
\newcommand{\out}{\Theta}

\newcommand{\nlayer}{k}
\newcommand{\nlayers}{K}
\newcommand{\nheads}{K}

%Edge and Node features
\newcommand{\ef}{\mathrm{w}}
\newcommand{\efhat}{\mathrm{\hat w}}
\newcommand{\nf}{\mathrm{f}}
\newcommand{\EF}{\mE}
\newcommand{\EFhat}{\hat\mE}

\newcommand{\nphis}{T}
\newcommand{\vtx}{v}

\newcommand{\absbound}{A}

\newcommand{\consteqfh}{C}
\newcommand{\constslrrls}{C}

\newcommand{\constsphiphis}{C}
\newcommand{\peqf}{\phi}
\newcommand{\tpeqfout}{\tilde{\phi}_{out}}
\newcommand{\peqfout}{\phi_{out}}
\newcommand{\peqfin}{\phi_{in}}

\newcommand{\HSdim}{{d_{h}}}
\newcommand{\PEdim}{{d_{p}}}
\newcommand{\mdim}{{d_{v}}}
\newcommand{\efdim}{{d_{e}}}
\newcommand{\nfdim}{{d_{n}}}
\newcommand{\outdim}{{d_{out}}}
\newcommand{\const}{C}


\newcommand{\constone}{C_1}
\newcommand{\consttwo}{C_2}


\DeclareMathOperator*{\Agg}{Agg}

\DeclareMathOperator*{\diag}{diag}

\newcommand{\contgraphspace}{\mathcal{G}^*}
\newcommand{\graphspace}{\mathcal{G}}
\newcommand{\edgesignalspace}{\mathcal{E}}
\newcommand{\nodesignalspace}{\mathcal{X}}




\newcommand{\adj}{a}
\newcommand{\maxval}[1]{M_{#1}}


\DeclareMathOperator*{\concat}{\mathbin\Vert}
% \newcommand{\BigDbar}{\mathop{\vcenter{\hbox{\Huge$\Vert$}}}\limits}


\newcommand{\regions}{\mathcal{J}}
\newcommand{\region}{\bm{I}}
\newcommand{\nsamples}{n}

\newcommand{\ptset}{\mathcal{S}}
% \newcommand{\ptsetspace}{\mathcal{X}}
% \newcommand{\ptsetspacehat}{\hat{\mathcal{X}}}

% \newcommand{\ptsetspacehat}{\hat{\mathcal{X}}}

% notation
\newcommand{\Space}{\mathcal{X}}
\newcommand{\token}{token\xspace}
\newcommand{\tokens}{tokens\xspace}
\newcommand{\tokenset}{tokenset\xspace}
\newcommand{\tokensets}{tokensets\xspace}

\newcommand{\Tokenset}{Tokenset\xspace}

\newcommand{\tocite}[1]{\textcolor{blue}{cite: #1}}


\newcommand{\hyp}[6]{\mathcal{H}^{(#1)}({#2, #3, #4, #5, #6})}
\newcommand{\hypothesis}{\bm{\mathcal{H}}}



	\newcommand{\dataset}{\mathcal{D}}
	\newcommand{\loss}{\mathcal{L}}

	\newcommand{\natdist}{P_{\geq N}}

	\newcommand{\sample}[1]{\sim_{#1}}


	\newcommand{\outEE}{
		\mathop{\mathbb{E}}\limits_{%
			\substack{
				% \{(\mathcal{T}_i,\,y_i)\}\sim \nu^m 
			% \\
			n_i \sim \natdist
		}}
	}

% template token set
\newcommand{\TT}{\bm{\mathcal{T}}}
% sampled token set
\newcommand{\ST}{\bm{T}}
\newcommand{\admconsts}{(L_{\nf}, C_\chi, D_\chi)}


\newcommand{\outE}{
		\mathop{\mathbb{E}}\limits_{%
			 \substack{
				\{\TT_i,\,y_i\}_{i=1}^m \sim \nu^m 
			\\
			n_i \sim \natdist
			\\
			{\ST}_i\sample{n_i} \TT}
		}
	}
	% \newcommand{\stdef}{(\sampledset, \nf, \hat\ef)}
\newcommand{\stdef}{(\sampledset, \nf)}
% \newcommand{\stdefi}[1]{(\sampledset_{#1}, \nf_{#1}, \hat\ef_{#1})}
\newcommand{\stdefi}[1]{(\sampledset_{#1}, \nf_{#1})}

\newcommand{\ttdef}{((\chi, \dist, \mu), \nf)}
% \newcommand{\ttdef}{((\chi, \dist, \mu), \nf, \ef)}
\newcommand{\ttdefi}[1]{((\chi_{#1}, \dist_{#1}, \mu_{#1}), \nf_{#1})}
% \newcommand{\ttdefi}[1]{((\chi_{#1}, \dist_{#1}, \mu_{#1}), \nf_{#1}, \ef_{#1})}
\newcommand{\eventtTS}[3]{\mathcal{E}(#1, #2, #3)}
\newcommand{\eventtRS}[3]{\mathcal{E}_{R}(#1, #2, #3)}
\newcommand{\eventT}[3]{\mathcal{E}(#1, #2, #3)}

\newcommand{\emp}{\mathrm{emp}}
\newcommand{\RNum}[1]{\uppercase\expandafter{\romannumeral #1\relax}}
\newcommand{\Sconst}{\| \msg \|_{\infty} \sqrt{C_{\chi}} \tau}

\newcommand{\constI}{C_{I}}
\newcommand{\constIdef}{4 (\Lip{\att} \|\msg\|_{2, \infty} + \Lip{\msg})}
\newcommand{\constII}{C_{II}}
\newcommand{\constIIdef}{ \left(4\LLip{\att}  \LLip{\PE} \maxval{\msg} + \LLip{\msg}  \LLip{\ef} + \maxval{\msg} \LLip{\ef}\right)}
\newcommand{\constIII}{C_{III}}
\newcommand{\tauUB}{\tfrac{1}{4\sqrt{C_\chi}} n^{\nicefrac{1}{(D_\chi + 2)}}}

\newcommand{\PEstable}{\rho}
\newcommand{\PEparams}{(\rho, R)}
\newcommand{\universe}{\mathcal{U}}

\newcommand{\offdiag}{\mathrm{offdiag}}

\newcommand{\roughconst}{54}






\newcommand{\ncondition}[1]{#1^{\frac{1}{D_\chi + 2}} \geq 24C_\chi\log #1}


\newcommand{\normadj}{p}
\newcommand{\mPi}{\bm{\Pi}}
\newcommand{\mDelta}{\bm{\Delta}}



\newcommand{\HIdef}{16^{\nlayers + 1} L_{\pool}  \nlayers \LLip{\msg}^{\frac{\nlayers(\nlayers + 1)}{2}}
\LLip{\att}^{\nlayers} 
\left(\LLip{\nf} + \nlayers \LLip{\PE}\right)  \sqrt{C_\chi}}

\newcommand{\HIdefO}{\Lip{\msg}  \Lip{\att} + \sqrt{C_\chi} \tau}
\newcommand{\HIIdef}{16^{\nlayers + 1} L_{\pool}  \nlayers \LLip{\msg}^{\frac{\nlayers(\nlayers + 1)}{2}}
					% \|\nf\|^{\nlayers}_{2, \infty} 
					\LLip{\att}^{\nlayers} 
					{C_\chi} }
\newcommand{\HIIdefO}{\Lip{\msg} }
\newcommand{\HIIIdef}{16^{\nlayers + 1} L_{\pool}  \nlayers \LLip{\msg}^{\frac{\nlayers(\nlayers + 1)}{2}}
					% \|\nf\|^{\nlayers}_{2, \infty} 
					\LLip{\att}^{\nlayers} R \tau }

					

\addbibresource{ref.bib}
\title{Learning  a \sout{function on a} manifold}
\date{December 2025}

\begin{document}
\maketitle

\newcommand{\manifold}{\mathcal{M}}
\renewcommand{\dataset}{{D}}


\section{Definitions}

\begin{definition}
    A metric tensor on (the embedded in $\mathbb{R}^k$) manifold $\manifold$ is a function $g: \manifold \to \mathbb{R}^{k \times k},\forall x \in \manifold, \rank g(x) = \dim \manifold, g(x) \succeq 0$.
\end{definition}


\begin{definition}[Dirichlet energy]
The \emph{Dirichlet energy} of a function $f : \manifold \to \mathbb{R}$ with respect to the manifold $\manifold$ equipped with the metric tensor $g$ is defined as
\[
    \mathcal{E}_{\manifold, g}(f) = \frac{1}{2} \int_{\manifold} \|\nabla_g f(x)\|_g^2 \, d\mu_g(x),
\]
where $\nabla_g f$ denotes the gradient of $f$ with respect to the metric $g$, $\|\cdot\|_g$ denotes the norm induced by $g$, and $\mu_g$ is the volume measure associated to $g$ on $\manifold$.
\end{definition}

\begin{definition}[Measure]
    A Riemannian metric on a $d$-dimensional manifold is a smoothly varying inner product
    $$
    g_x: T_x M \times T_x M \rightarrow \mathbb{R} .
    $$
    In local coordinates, it is represented by a positive-definite matrix
    $
    g(x)=\left[g_{i j}(x)\right] .
    $
    The metric induces the Riemannian volume measure, often denoted:
    $$
    d \mu_g(x)=\sqrt{\operatorname{det} g(x)} d x^1 \cdots d x^d .
    $$
    This is the analogue of ``area $=\sqrt{\operatorname{det} g}$'' in differential geometry.
    $d \mu_g(x)=\sqrt{\operatorname{det} g(x)} d x^1 \cdots d x^d$.
\end{definition}


\begin{definition}[Gradient]
    The gradient of a function $f: \manifold \to \mathbb{R}$ with respect to the metric tensor $g$ is defined as
    \[
        \nabla_g f(x) = g(x) \nabla f(x) \in T_x \manifold,
    \]
    where $\nabla f(x)$ is the gradient of $f$ with respect to the Euclidean metric.
\end{definition}

\begin{definition}[Laplacian]
    The Laplacian of a function $f: \manifold \to \mathbb{R}$ with respect to the metric $g$ is defined as
    \[
        \Delta_g f(x) = \nabla \cdot \nabla_g f(x) = \nabla \cdot (g \nabla f(x)) \in \mathbb{R}
    \]
\end{definition}

\section{Learning a function}
\paragraph{Given:} 
\begin{itemize}
    \item An embedded manifold $\manifold \subset \mathbb{R}^k, \dim \manifold = d$
    \item A metric tensor on $\manifold$: $g: \manifold \to \mathbb{R}^{k \times k},\forall \vx \in \manifold, \rank g(\vx) = d, g(\vx) \succeq 0$. 
    
    \emph{Note: For simplicity, we assume $\det g(\vx) = 1$ hence, the measure is the Lebesgue measure.}
    \item A dataset $\dataset = \{(\vx_i, y_i)\}_{i=1}^n$ where $\vx_i \in \manifold, y_i \in \mathbb{R}$ 
\end{itemize}

\paragraph{Goal:} to learn the smoothest function $f: \manifold \to \mathbb{R}$ such that the function $f(\vx_i) \approx y_i$. 

Constrained optimization:
\begin{align*}
    \min_{f} \mathcal{E}_{\manifold, g}(f) \\
    \text{s.t. } f(\vx_i) = y_i, \forall i 
     \in [n].
\end{align*}

Unconstrained optimization:
\begin{gather*}
    \min_{f} \mathcal{E}_{\manifold, g}(f) + \lambda \sum_{i=1}^n (f(\vx_i) - y_i)^2 
    = \\
    \min_{f}
    \frac{1}{2} 
    \int_{\manifold} \|\nabla_g f(\vx)\|_g^2 \, d\mu_g(\vx) + \lambda \sum_{i=1}^n (f(\vx_i) - y_i)^2
\end{gather*}
Gradient descent, convex optimization (in a \emph{function} space):
\begin{align*}
    \text{PDE: }& \frac{\partial f}{\partial t} = - \Delta_g f + 2\lambda \sum_{i=1}^n (f(\vx_i) - y_i) \delta_{\vx_i} 
    \\
    \text{step: }& f_{t+1} = f_t +  \eta\left(\Delta_g f_t - 2\lambda \sum_{i=1}^n (f(\vx_i) - y_i) \delta_{\vx_i}\right)
\end{align*}

\begin{definition}[Gradient-flow evolution operator]
    We will denote $\Phi_T(f_i)$ as the gradient-flow evolution operator of $f_i$ after time $T$ of gradient descent. Note, that $\Phi_T(f_i) = \Phi_T(f_i, g, \manifold, \dataset)$ depends on the metric tensor $g$, the manifold $\manifold$, and the dataset $\dataset = \{(\vx_i, y_i)\}_{i=1}^n$.
\end{definition}

\paragraph{Discretization in space}
(Disclaimer: this is a bit of a stretch.)
If $\manifold$ is discretized as a point cloud $\mathcal{P} = \{\vz_i\}_{i=1}^n$ with weights 
$$w_{ij} = \exp\left(-\frac{\|\vz_i - \vz_j\|_{g(\vz_i)}^2}{\epsilon}\right)$$
(epsilon TBD), one may argue that:
\begin{align*}
    L_{\mathcal{P}} D^{-1}_{\mathcal{P}}  \approx  \Delta_g 
\end{align*}
Hence, one step of gradient descent becomes (yes, here comes the "stretch") like a transformer layer:
\begin{align*}
    f_{t+1}(\vz) 
    = 
    \underbrace{f_t(\vz)}_{\text{Residual}} 
    +  
    \eta\left( 
    \underbrace{\frac{\sum_{i} \exp(-{\|\vz - \vz_i\|_{g(\vz)}^2}/{\epsilon}) f_t(\vz_i)}{\sum_{j} \exp(-{\|\vz - \vz_j\|_{g(\vz)}^2}/{\epsilon})}}_{\approx \text{Attention}}  - \underbrace{2\lambda \sum_{i=1}^n (f(\vx_i) - y_i) \delta_{\vx_i}(\vz)}_{\text{Residual}}
    \right) 
\end{align*}

\textbf{Idea 1: what if a fixed transformer performs gradient descent on the function space, hence, learns a function in context?} 


\section{Learning a manifold}

What if we are given a few smooth functions $f_1, \ldots, f_l: \mathbb{R}^k \to \mathbb{R}$ on a manifold $\manifold$. However, we don't know the metric tensor $g$ on $\manifold$. Can we learn the metric tensor such that the functions $f_1, \ldots, f_l$ are as smooth as possible? If we do so, then we can easily learn a new smooth function from a few observations on the fly next time. 

For this problem, the objective is the following:
\begin{gather*}
    \min_{g} \sum_{i=1}^l \mathcal{E}_{\manifold, g}(f_i) \\
    \text{s.t. } \det g(\vx) = 1 \text{ for all } \vx \in \manifold.
\end{gather*}
where $\mathcal{E}_{\manifold, g}(f_i)$ is the Dirichlet energy of $f_i$ with respect to the metric tensor $g$. If we discard the constraint on the determinant or relax it to $\alpha I \preceq g(\vx)^{-1} \preceq \beta I$, the problem is convex w.r.t. $g$.


% Given the fixed metric tensor $g$, we can find the smooth proxy after performing $T$ steps of gradient descent on the function space for each function $f_i$ if . Hence, the problem becomes:
% \begin{gather*}
%     \min_{\substack{g \\ g(\cdot) \succeq 0 \\ \det g(\cdot) = 1}} \sum_{i=1}^l \mathcal{E}_{\manifold, g}(\Phi_T(f_i, g, \manifold, \dataset))
% \end{gather*}


\textbf{Idea 2: training a transformer-like model = learning a geometry of a manifold} 

Okay, as of now, we didn't have a need to perform GD in a function space inside the global optimization with respect to the metric tensor $g$ -- there was no need to train a transformer-like model with multiple layers. In what case would that be necessary? 


\section{More Complex Geometry}

Okay, let us consider a more complex geometry so that performing gradient descent in a function space is necessary during optimization over $g$. 

What if the metric tensor $g(x)$ also depends on the function $f(x)$ on the manifold: $g(x, f)$? 

\paragraph{The Perona--Malik anisotropic diffusion} Equation is given by
\[
\frac{\partial u}{\partial t}
= \nabla \cdot \bigl( c(|\nabla u|) \, \nabla u \bigr),
\]
where \(u(x,t)\) is the evolving signal or image, and the 
diffusivity function \(c(\cdot)\) is typically chosen as
\[
c(s) = \exp\!\left( - \left(\frac{s}{K}\right)^2 \right)
\qquad\text{or}\qquad
c(s) = \frac{1}{1 + \left(\frac{s}{K}\right)^2 }.
\]

\section*{Relevant literature}

\begin{itemize}
    \item \textcite{geshkovskiMathematicalPerspectiveTransformers2025} -- models transformers as interacting particle systems
    \item \textcite{belkinConvergenceLaplacianEigenmaps2006} -- on the convergence of discrete Laplacian to Laplace-Beltrami operator
    \item \textcite{zhangPDETransformerContinuousDynamical2025} -- PDE view on transformers
\end{itemize}

\printbibliography

\end{document}